% ============================================================
% ANNEXE B : CODE SOURCE COMPLET
% ============================================================

\chapter{Extraits de Code Source Complémentaires}

\section{Script Principal (main.py)}

\begin{lstlisting}[style=pythonstyle, caption={Point d'entrée principal du système DMS.}, label={lst:main}]
#!/usr/bin/env python3
"""
DMS - Driver Monitoring System
Point d'entree principal.

Usage:
    python main.py                    # Mode simulation
    python main.py --scenario mixed   # Scenario specifique
    python main.py --mode robot       # Mode prototype robot
    python main.py --mode camera      # Mode camera en direct
"""

import argparse
import sys
import time
from src.video_acquisition import VideoAcquisition
from src.face_eye_detection import FaceEyeDetector
from src.head_pose import HeadPoseEstimator
from src.behavioral_analysis import BehavioralAnalyzer
from src.ecu_decision import ECUDecisionModule
from src.scenarios import ScenarioRunner
from src.visualization import Visualizer


def parse_args():
    parser = argparse.ArgumentParser(
        description='DMS - Driver Monitoring System'
    )
    parser.add_argument('--mode', type=str, 
        default='simulation',
        choices=['simulation', 'camera', 'robot'],
        help='Mode de fonctionnement')
    parser.add_argument('--scenario', type=str,
        default='all',
        choices=['attentive', 'fatigued', 'distracted', 
                 'mixed', 'all'],
        help='Scenario de test')
    parser.add_argument('--duration', type=int, 
        default=120,
        help='Duree de simulation (secondes)')
    parser.add_argument('--fps', type=int, default=30,
        help='Images par seconde')
    parser.add_argument('--camera', type=int, default=0,
        help='Index de la camera')
    parser.add_argument('--output', type=str, 
        default='results',
        help='Dossier de sortie des resultats')
    return parser.parse_args()


def run_simulation(args):
    """Execute les scenarios de simulation."""
    print("=" * 60)
    print("  DMS - Driver Monitoring System")
    print("  Mode: Simulation")
    print("=" * 60)
    
    runner = ScenarioRunner(
        duration=args.duration,
        fps=args.fps
    )
    
    if args.scenario == 'all':
        scenarios = ['attentive', 'fatigued', 
                     'distracted', 'mixed']
    else:
        scenarios = [args.scenario]
    
    results = {}
    for scenario in scenarios:
        print(f"\n--- Scenario: {scenario} ---")
        result = runner.run(scenario)
        results[scenario] = result
        
        # Affichage des metriques
        print(f"  Score fatigue moyen: "
              f"{result['avg_fatigue']:.3f}")
        print(f"  Score fatigue max: "
              f"{result['max_fatigue']:.3f}")
        print(f"  Score distraction moyen: "
              f"{result['avg_distraction']:.3f}")
        print(f"  PERCLOS moyen: "
              f"{result['avg_perclos']:.1f}%")
        print(f"  Alertes generees: "
              f"{result['total_alerts']}")
    
    # Visualisation
    viz = Visualizer(output_dir=args.output)
    viz.plot_all_scenarios(results)
    
    print("\n" + "=" * 60)
    print("  Simulation terminee avec succes!")
    print(f"  Resultats sauvegardes dans: {args.output}/")
    print("=" * 60)


def run_camera(args):
    """Execute le DMS avec une camera en direct."""
    print("DMS - Mode Camera en Direct")
    
    # Initialisation des modules
    video = VideoAcquisition(
        source=args.camera, fps=args.fps
    )
    detector = FaceEyeDetector()
    pose_estimator = HeadPoseEstimator()
    analyzer = BehavioralAnalyzer(fps=args.fps)
    ecu = ECUDecisionModule()
    
    if not video.open():
        print("Erreur: impossible d'ouvrir la camera")
        sys.exit(1)
    
    try:
        for frame, gray in video.generate_frames():
            # Detection
            detections = detector.detect(frame, gray)
            
            if detections:
                det = detections[0]
                eye_state = det['eye_state']
                
                # Pose de la tete (si landmarks)
                head_pose = {'yaw': 0, 'pitch': 0, 
                            'roll': 0}
                
                # Analyse comportementale
                analysis = analyzer.update(
                    eye_state, head_pose
                )
                
                # Decision ECU
                decision = ecu.update(
                    analysis['fatigue_score'],
                    analysis['distraction_score']
                )
                
                # Affichage
                print(f"\rEtat: {analysis['driver_state']} "
                      f"| Fatigue: "
                      f"{analysis['fatigue_score']:.2f} "
                      f"| Alerte: "
                      f"{decision['alert_level']}", 
                      end='')
    
    except KeyboardInterrupt:
        print("\nArret du systeme.")
    finally:
        video.release()


if __name__ == '__main__':
    args = parse_args()
    
    if args.mode == 'simulation':
        run_simulation(args)
    elif args.mode == 'camera':
        run_camera(args)
    elif args.mode == 'robot':
        # Import conditionnel pour le mode robot
        from src.robot_controller import RobotDMS
        robot = RobotDMS(camera=args.camera, fps=args.fps)
        robot.run()
\end{lstlisting}

\section{Module de Scénarios (scenarios.py)}

\begin{lstlisting}[style=pythonstyle, caption={Extrait du module de scénarios de test.}, label={lst:scenarios}]
class ScenarioRunner:
    """Executeur de scenarios de test DMS."""
    
    def __init__(self, duration=120, fps=30):
        self.duration = duration
        self.fps = fps
        self.analyzer = BehavioralAnalyzer(fps=fps)
        self.ecu = ECUDecisionModule()
    
    def run(self, scenario_name):
        """Execute un scenario de test."""
        # Generation des signaux
        signals = self._generate_signals(scenario_name)
        
        # Execution
        results = {
            'fatigue_scores': [],
            'distraction_scores': [],
            'perclos_values': [],
            'alert_levels': [],
            'driver_states': []
        }
        
        total_frames = self.duration * self.fps
        
        for i in range(total_frames):
            t = i / self.fps
            
            # Etat des yeux et pose
            eye_state = signals['eye_states'][i]
            head_pose = {
                'yaw': signals['yaw'][i],
                'pitch': signals['pitch'][i],
                'roll': 0.0
            }
            
            # Analyse
            analysis = self.analyzer.update(
                eye_state, head_pose, t
            )
            
            # Decision ECU
            decision = self.ecu.update(
                analysis['fatigue_score'],
                analysis['distraction_score'], t
            )
            
            # Enregistrement
            results['fatigue_scores'].append(
                analysis['fatigue_score'])
            results['distraction_scores'].append(
                analysis['distraction_score'])
            results['perclos_values'].append(
                analysis['perclos'])
            results['alert_levels'].append(
                decision['alert_level'])
            results['driver_states'].append(
                analysis['driver_state'])
        
        # Metriques resumees
        results['avg_fatigue'] = np.mean(
            results['fatigue_scores'])
        results['max_fatigue'] = np.max(
            results['fatigue_scores'])
        results['avg_distraction'] = np.mean(
            results['distraction_scores'])
        results['avg_perclos'] = np.mean(
            results['perclos_values'])
        results['total_alerts'] = sum(
            1 for a in results['alert_levels'] if a > 0)
        
        return results
\end{lstlisting}

\section{Module de Visualisation (visualization.py)}

\begin{lstlisting}[style=pythonstyle, caption={Extrait du module de visualisation.}, label={lst:visualization}]
import matplotlib.pyplot as plt
import numpy as np

class Visualizer:
    """Visualisation des resultats DMS."""
    
    def __init__(self, output_dir='results'):
        self.output_dir = output_dir
        os.makedirs(output_dir, exist_ok=True)
    
    def plot_scenario(self, name, results, duration):
        """Genere les graphiques pour un scenario."""
        t = np.linspace(0, duration, 
                        len(results['fatigue_scores']))
        
        fig, axes = plt.subplots(4, 1, figsize=(14, 10),
                                 sharex=True)
        fig.suptitle(f'Resultats DMS - {name}', 
                     fontsize=14, fontweight='bold')
        
        # PERCLOS
        axes[0].plot(t, results['perclos_values'], 
                     'b-', linewidth=1.5)
        axes[0].axhline(y=20, color='g', linestyle='--',
                        label='Normal')
        axes[0].axhline(y=40, color='y', linestyle='--',
                        label='Fatigue')
        axes[0].axhline(y=70, color='r', linestyle='--',
                        label='Somnolence')
        axes[0].set_ylabel('PERCLOS (%)')
        axes[0].legend()
        axes[0].grid(True)
        
        # Score de fatigue
        axes[1].plot(t, results['fatigue_scores'], 
                     'r-', linewidth=1.5)
        axes[1].axhline(y=0.35, color='y', linestyle='--')
        axes[1].axhline(y=0.6, color='orange', 
                        linestyle='--')
        axes[1].axhline(y=0.75, color='r', linestyle='--')
        axes[1].set_ylabel('Score Fatigue')
        axes[1].grid(True)
        
        # Score de distraction
        axes[2].plot(t, results['distraction_scores'], 
                     'm-', linewidth=1.5)
        axes[2].set_ylabel('Score Distraction')
        axes[2].grid(True)
        
        # Niveau d'alerte
        axes[3].step(t, results['alert_levels'], 
                     'k-', linewidth=2)
        axes[3].set_ylabel('Niveau Alerte')
        axes[3].set_xlabel('Temps (s)')
        axes[3].set_yticks([0, 2, 3, 4])
        axes[3].set_yticklabels(['Normal', 'Warning', 
                                 'Critical', 'Emergency'])
        axes[3].grid(True)
        
        plt.tight_layout()
        filepath = os.path.join(self.output_dir, 
                               f'{name}.png')
        plt.savefig(filepath, dpi=150)
        plt.close()
\end{lstlisting}
